





\documentclass[10pt, a4paper, titlepage]{article}

\usepackage[T1]{fontenc}
\usepackage[utf8]{inputenc}
\usepackage{lmodern}
\usepackage{textcomp}
\usepackage{microtype}
\usepackage[english]{babel}
\usepackage{csquotes}
% Bibliographie -----------------------------------------
\usepackage[backend=biber,
			style=alphabetic,
			isbn=false,
			doi=false
			]{biblatex}
\addbibresource{./content/references.bib}

% Typographie -------------------------------------------
\renewcommand{\theenumi}{(\roman{enumi})}
\usepackage{graphicx}
\usepackage{setspace}
% -------------------------------------------------------

% Mathe Pakete, Einstellungen und Commands --------------
\usepackage{amsmath}
\usepackage{amsfonts}
\usepackage{amssymb}
\usepackage{amsthm}
\usepackage{euscript}
\numberwithin{equation}{section}


% Theorem-Umgebungen ------------------------------------
\usepackage{thmtools}

\declaretheorem[
	name=Theorem,
	numberwithin=section
	]{theorem}
\declaretheorem[
	name=Lemma,
	sibling=theorem,
	]{lemma}
\declaretheorem[
	name=Proposition,
	sibling=theorem,
	]{prop}
\declaretheorem[
	name=Corollary,
	sibling=theorem,
	]{corollary}

\declaretheorem[
	name=Definition,
	style=definition,
	numbered=no,
	]{definition}

\declaretheorem[
	name=Remark,
	style=remark,
	numbered=no
	]{remark}	
\declaretheorem[
	name=Example,
	style=definition,
	numbered=no
	]{example}


% Hyperref (immer am Schluss der Präambel!) -------------
\usepackage[pdftex,pdfpagelabels]{hyperref}
\hypersetup{pdftitle={Beweis des Satz von Thales},
			pdfauthor={Maxime Musterfrau}}
% -------------------------------------------------------

%%%%%%%%%%%%%%%%%%%%%%%%%%%%%%%%%%%%%%%%%%%%%%%%%%%%%%%%%
%%% Beginn des Dokuments
%%%%%%%%%%%%%%%%%%%%%%%%%%%%%%%%%%%%%%%%%%%%%%%%%%%%%%%%%
\begin{document}

% Titelei -----------------------------------------------
\hypersetup{pageanchor=false}
%!TEX root = ../thesis.tex
\begin{titlepage}
\vspace*{-2cm}  % bei Verwendung von vmargin.sty
\begin{flushright}
    \includegraphics[width=8cm]{./content/Uni_Logo_2016_SW.jpg}
\end{flushright}
%\vspace{1cm}

\begin{center}  % Diplomarbeit ODER Magisterarbeit ODER Dissertation
    \Huge{\textbf{\textsf{\MakeUppercase{
        Bachelorarbeit
    }}}}
    \vspace{2cm}

    \large{\textsf{  % Diplomarbeit ODER Magisterarbeit ODER Dissertation
                     % (Dies ist erst die Ueberschrift!)
        Titel der Bachelorarbeit
    }}
    \vspace{.1cm}

    \LARGE{\textsf{  Analytic sets
    }}
    \vspace{3cm}

    \large{\textsf{  % Verfasserin ODER Verfasser (Ueberschrift)
        Verfasser
    }}

    \Large{\textsf{  Charlie Ahrendts
    }}
    \vspace{3cm}

    \large{\textsf{
        angestrebter akademischer Grad  % (Ueberschrift)
    }}

    \Large{\textsf{  % Magistra ODER Magister ODER Doktorin ODER Doktor
                     % ACHTUNG: Kuerzel "Mag.a" oder "Dr.in" nicht zulaessig
        Bachelor of Science (BSc.)
    }}
\end{center}
\vspace{2cm}

\noindent\textsf{Wien, im Monat September 2025}  % <<<<< ORT, MONAT UND JAHR EINTRAGEN
\vfill

\noindent\begin{tabular}{@{}ll}
\textsf{Studienkennzahl lt.\ Studienblatt:}
&
\textsf{UA 033621}  % <<<<< STUDIENKENNZAHL EINTRAGEN
\\
    % BEI DISSERTATIONEN:
%\textsf{Dissertationsgebiet lt. Studienblatt:}
    % ANSONSTEN:
\textsf{Studienrichtung lt.\ Studienblatt:}
&
\textsf{Mathematik}  % <<<<< DISSGEBIET/STUDIENRICHTUNG EINTRAGEN
\\
% Betreuerin ODER Betreuer:
\textsf{Betreuer: }
&
\textsf{Dipl.-Phys. Dipl.-Math. Dr. Peter Elbau}  % <<<<< NAME EINTRAGEN
\end{tabular}

\end{titlepage}


\newpage

\thispagestyle{empty}

%!TEX root = ../thesis.tex
\section*{Abstract}

Analytic sets form an often overlooked but nontheless very useful field of study. 
Their history teaches us about how mistakes lead to new discoveries and their current uses show how wanting to prove someone wrong can lead to important results in many different areas of mathematics. 

Analytic sets will be introduced as continuous images of Polish spaces. 
We will work towards an understanding of their nature, especially concerning differences and commonalities with Borel sets.

A main objective of this thesis will be a  proof of the fact that there exist Analytic sets which are not Borel. 
Further time will be spent on the measurability properties of Analytic sets and different frameworks in which they appear.



\newpage

%\thispagestyle{empty}
\tableofcontents

\newpage
% -------------------------------------------------------
\hypersetup{pageanchor=true}
\renewcommand{\thepage}{ \arabic{page} }

\setcounter{page}{1}
\onehalfspacing







\section{A short history lesson}
As it is so often the case with new and interesting discoveries, mathematical areas of study often come about unplanned or on accident. 
One prime example of this is the field of Analytic sets. Studying it provides us, next to some important results in among others set theory, probability theory and functional analysis with a deeper understanding of the way mathematics is done in practice.\\
Our story starts in the Year 1905 with famous Frensch mathematician Henri Lebesgue and his paper "Sur les fonctions représentables analytiquement" \cite{lebesgue1905}. 
In it, Lebesgue proved many statements of great importance. 
He did however include one small remark which stated, without proof, that the Borel sets are closed under Projections. \\
About 10 years later, Russian mathematician Nikolai Luzin ordered his student Souslin to study Lebesgues paper. 
When the still very young Souslin returned, saying he had found an error in Lebesgues work, Luzin and his colleage Sierpiński to their credit believed their student. \cite{sierpinski1950}
Soon after, they would be the ones leading the charge on this young new field of mathematics that had just opened itself up to them. 
As Souslin had correctly identified, there was no justification of the claim that projections of Borel sets would again be Borel. \\
In a short publication in the French Journal "Comptes Rendu" in 1917, Souslin annouced the falsity of Lebesgues statement and gave a rough idea of his proof. [cite sur une ]
These new sets that arise out of projections of Borel sets are called Analytic sets or Souslin sets.\\
Much more thorough treatments of the topic as well as the first construction of an analytic set that is not Borel were later given by Luzin and Sierpinski after Souslins early death in 19 [cite sur quelques proprietes\ldots], \cite{lusin1923} \cite{lusin1927}
It took until after the end of World War I for western mathematicians to start engaging with this new field. But once they did, it quickly started to become apparent that Analytic sets were going to become a very important tool for many different areas of mathematics. \cite{Moschovakis1987} \cite{rogers1980}.\\
This thesis aims to give an overview over some of the basic concepts, theorems and proofs in the field of Analytic sets. Through chapters 2-4, we will mainly follow the conventions and proofs of \cite{cohn2013}



%%%%%%%%%%%%%%%%%%%%%%%%%%%%%%%%%%%%%%%%%%%%%%%%%%%%%%%%%%%%%

\section{Polish spaces and Analytic sets}
As previously stated, Analytic sets arose out of the question, whether projections of Borel sets are again Borel and Souslins counterproof of that claim. 
So if we seek to study the greater class of sets that is obtained as projections of Borel sets, it would only make sense do define them exactly as such.\\ 
As it turns out however, there are many different ways to define these sets, all of them useful in their own ways. 
The one that will be used throughout the following chapters is one that allows for particularly nice versions of many of the basic results considered here.
We will define Analytic sets as continuous images of Polish spaces, named in honor of the Polish mathematicians who were the first to extensively study them. [cite]



\begin{definition}[Polish space]
	A topological space $X$ is called a Polish space, if it is completely metrizable and seperable (contains a countable dense subset)
\end{definition}

Interesting to note here is the difference between a complete metric space and a completely metrizable space.
For the latter, we only require the existence of a complete metric on $X$, but we do not need to choose a concrete one.
This means of course that all complete metric spaces are Polish\\
\\
There are a few different ways to view analytic sets, as well as different ways to define them. 
Souslin, for example in his early works on the topic defined them as arising from a series of unions and intersections of certain families of sets.
We shall here however stick to the most common, and often most useful  definition, which is that of Analytic sets being continuous images of Polish spaces 

\begin{definition}[Analytic set]
	Let $X$ be a polish space, $A \subset X$. We call A analytic, if there exists a Polish space $Y$ and  $f:Y \to X$ continuous, such that  $f(Y) = A$

\end{definition}

Throughout this section, we will use some standard results about topological and metric spaces without proof. These are:
\begin{corollary}
	\label{second_countable}
	In metric spaces, seperability and second countablility are equivalent
\end{corollary}
\begin{corollary}
	\label{closed_cc}
	closed subsets of complete metric spaces are again complete metric spaces
\end{corollary}
Proofs for these theorems can be found in [cite/Anhang]

\begin{lemma}
	\label{products}
	Finite and countable products of Polish spaces are polish.
\end{lemma}
\begin{proof}
	Let $X_1,X_2,\ldots$ be a sequence of (nonempty) Polish spaces. We can choose a complete metric $\overline{d_i}$ for each of the spaces $X_i$. 
	Now for each $i$, let $d_i \left( x,y \right) := min \left\{ 1,\overline{d_i} \left( x,y \right)  \right\}$\\
	This again defines a complete metric with the additional property that $d_i\left( x,y \right) \leq 1$ for all $x,y \in X_i$ [Genauer beweisen? Anhang?]
This new metric retains only information about small distances in $\overline{d_i}$.\\
Now we can turn towards the cartesian product $X := \prod X_i$.
Let $d\left( x,y \right):= \sum \frac{1}{2^{n}} d_i\left( x_{i},y_{i} \right)$
where $x = \left( x_1,x_2,\ldots \right), y = \left( y_1,y_2,\ldots \right) \in X$ 
This sum converges for all $x,y \in X$ since we constructed the $d_i$ to be bounded by 1.
 Moreover, $d$ defines a metric on $X$. This is easily seen by the fact that positive definieteness, symmetry and the triangle inequality all hold in each term of our sum individually and thusfor the whole sum. 
 The topology generated by $d$ is exactly the product topology on $X$ [Genauer?].\\

[Ab hier nicht im Buch, hoffe der Beweis passt so] \\
To show that $d$ is indeed complete, we take an arbitrary Cauchy sequence $(x_{n})_{n \in  \mathbb{N}}$ in $X$. The projection onto each  space $X_i$ is also a cauchy sequence  $x_{i_n}$.
 Since $X_i$ is complete, $x_{i_n}$ converges to some value $x_i$. 
Let $x:= \left( x_1,x_2,\ldots \right) $ be the componentwise limit of our Cauchy sequences. We need to show that $x_{n}$ converges to $x$ in the metric $d$:\\
Choose an arbitrary $\varepsilon > 0 $ and $m \in  \mathbb{N}$, such that $\sum_{i \geq m} \frac{1}{2^i} \leq \frac{\varepsilon}{2} $.
For each $i \leq m$, we can find an  $N_i$, such that $d_i\left( x_{i_n},x_i \right) \leq \frac{\varepsilon}{2}$ for all $n\geq N_i$. 
If we choose $N := max_{i \leq m}\left\{ N_i \right\} $, we get the following esimate for $n \geq N$: 

\begin{align*}
	d\left( x_{n},x \right) &= \sum \frac{1}{2^i} d_i\left( x_{i_n},x_i \right)\\
				&= \sum_{i\leq m} \frac{1}{2^i} d_i\left( x_{i_n},x_i \right) + \sum_{i \geq m} \frac{1}{2^i} d_i\left( x_{i_n},x_i \right) \\
				&\leq \sum_{i \leq m} \frac{1}{2^i} \frac{\varepsilon}{2} + \sum_{i\geq m} \frac{1}{2^i} \\
				&\leq \frac{\varepsilon}{2} + \frac{\varepsilon}{2} = \varepsilon
\end{align*}
So we get convergence of $x_n$ and thus completeness of the metric space $\left( X,d \right) $ \\

What remains to be shown is seperability. We want to make use of the equivalence of seperability and second countability for metric spaces. [cite]
For each $i$, we can find a countable Basis  $\mathcal{U}_i$ for the topology on $X_i$.
The sets of the form $U_1 \times U_2 \times \ldots \times U_N \times X_{N+1} \times X_{N+2} \times \ldots$ form a countable Basis for the product topology on $X$, 
so  $X$ is a seperable metrisable space and thus Polish.

\end{proof}

\begin{corollary}
	Open and closed subsets of polish spaces are polish
\end{corollary}


\begin{proof}
	Let $X$ be a polish space, $A$ be an open or closed subset of $X$
	By Corollary \ref{second_countable}, $X$ has a countable basis of open sets for its Topology. Then the restrictions of those sets to  $A$ form  a countable basis for the subspace topology on $A$. 
By the same equivalence, we get seperability of  $A$.\\
If $A$ is closed, it is according to Corollary \ref{closed_cc} completely metrizable by any complete metric on $X$, resticted to  $A$.\\
Now only the case of  $A$ being an open subset of  $X$ remains.
Let  \[
d_0 \left( x,y \right) := d\left( x,y \right) + \left| \frac{1}{d\left( x,A^{C} \right) } - \frac{1}{d\left( y,A^C \right) } \right| 
.\]
where $d$ is a complete metric on  $X$ and  $d\left( x,A^C \right) := \inf\left( d\left( x,z \right): z \in A^C \right)  $.\\
$d_0$ defines a metric on $A$. [proof?]

\end{proof}


\begin{theorem}
	Open and closed subsets of Polish spaces are analytic
\end{theorem}

\begin{proof}
	Let $X$ be a Polish space, $A \subset X$ open or closed.
	We know that $A$, as a space is polish. So there exists a polish space $Y$ and a continuuous function  $f:Y\to A$, such that  $f(Y) = A$. 
	Let $\iota := A \hookrightarrow X$ be the canonical embedding of  $A$ into  $X$. 
	Then we can define  $\tilde{f} := f \circ \iota$. Then  $\tilde{f}:Y \to X $ and  $\tilde{f}(Y) = A$ hold. So  $A$ is an analytic subset of X


\end{proof}

Two particular spaces that are of great interest in the study of analytic sets are $\mathbb{N}^{\mathbb{N}}$ and $\left\{ 0,1 \right\}^{\mathbb{N}} $.
We will see that the polish space $Z$ in our definition of analytic set can always be replaced by the space $\mathbb{N}^{\mathbb{N}}$. But first we need to verify that they are in fact polish:

\begin{theorem}
	$\mathbb{N}^\mathbb{N}$ is polish
\end{theorem}
\begin{proof}
	by $\mathbb{N}$ we always mean the natural numbers together with the discrete topology in which every subset is open. Since $\mathbb{N}$ is countable, seperability immediately follows
	The discrete metric $d\left( m,n \right) = 1 - \delta_{mn} $, which equals is a complete metric on  $\mathbb{N}$.
	In this metric the only Cauchy sequences are those that are eventually constant which obviously converge.
	It now follows from \ref{products}, that the cartesian product $\mathbb{N}^{\mathbb{N}}$ is also Polish.
\end{proof}

\begin{theorem}
	$\left\{ 0,1 \right\}^\mathbb{N}$ is polish
\end{theorem}
\begin{proof}
	The proof for this is equivalent to that of $\mathbb{N}^{\mathbb{N}}$ being polish. We again choose the discrete topology and discrete metric on $\left\{ 0,1 \right\} $ and use \ref{products} to get that $\left\{ 0,1 \right\} ^{\mathbb{N}} $ is Polish.	
\end{proof}


\begin{theorem}
	Let $X$ be a Polish space. Then there is a continuous function $f: \mathbb{N}^\mathbb{N} \to X $, such that $f\left( \mathbb{N}^\mathbb{N} \right) = X$
\end{theorem}


\begin{theorem}
	Let $A$ be a nonempty analytic subset of a polish space $X$. Then there exists continuous function $f: \mathbb{N}^{\mathbb{N}} \to X$, such that $f\left( \mathbb{N}^{\mathbb{N}} \right) = A$
\end{theorem}


%%%%%%%%%%%%%%%%%%%%%%%%%%%%%%%%%%%%%%%%%%%%%%%%%%%%%%%%%%

\section{Borel-gedöns}
The sets maybe closest in nature to the analytic sets are the Borel sets. 
This of course is a natural consequence of the way analytic sets were discovered. 
In some sense, Analytic sets are a just a generalization of the Borel sets which are closed under projections. 
We will try to formalise this notion as well as some other results about these two classes of sets in this section.

\begin{definition}[Borel set]
	We call a subset of a topological space Borel, if it is a member of $\mathcal{B}$, the $\sigma-$Algebra generated by the open sets
\end{definition}


\begin{theorem}
	Let $B$ be a Borel subset of a Polish space X. Then B is analytic		
\end{theorem}


\begin{definition}[Zero-Dimensional space?]
	.	
\end{definition}

\begin{theorem}
	Let B be a Borel subset of a Polish space X. Then there exists a zero-dimensional space Z, such that $f(Z) = B$
\end{theorem}

\begin{theorem}
	Let $B$ be an uncountable Borel subset of a polish space  $X$. Then  $B$ contains a subset which is homeomorphic to  $\left\{ 0,1 \right\}^\mathbb{N}$
	
\end{theorem}


\begin{theorem}[Separation theorem]
	.
\end{theorem}


\begin{theorem}
	Let A be a subset of a polish space X. If $A$ and  $A^{C}$ are analytic, then $A$ ist Borel.  
\end{theorem}


\begin{theorem}
	There exists.
\end{theorem}

\begin{definition}[Borel isomorphic]
	We call two Borel subsets $A,B$ of a Polish space X Borel isomorphic, if there exists a bijective, Borel measurable function $f: A \to B$
\end{definition}


\begin{theorem}
	Two Borel subsets of a polish space X are Borel isomorphic iff they have the same cardinality
\end{theorem}

We have seen that all Borel sets are analytic, but have not yet said anything about the converse. So we turn to the foundational theorem, introduced by souslin, that makes this field of analytic sets worth studying:
\begin{definition}[Universal set]
	
\end{definition}

\begin{theorem}[Souslin]
	There exist Analytic sets which are not Borel sets. 
\end{theorem}
\begin{proof}
	[Kechris p.85]
\end{proof}



\begin{example}
	As almost all 'nice' sets are Borel, we can assume that most constructions of Analytic non-borelian sets are fairly complicated. 
	One of the earliest such examples was provided in 1936 by polish mathematician Stefan Mazurkiewicz [cite]:\\
	Let $\mathbb{R}^I$ denote the space of real-valued functions, which are continuous in the closed interval $I:=[0,1]$. Let  $\Gamma$ be the set of functions  $f \in \mathbb{R}^I$, which are differentiable in $I$. This set is Co-analytic, meaning it is the complement of an analytic set, but is itself not analytic. 
	Thus, its complement cannot be Borel (or $\Gamma$ would be Borel as well, and thus Analytic, leading to a contradiction). \\
	The proof of this is rather long and technical and beyond the scope of this thesis, but can be found in [cite]
	
\end{example}

%%%%%%%%%%%%%%%%%%%%%%%%%%%%%%%%%%%%%%%%%%%%%%%%%%%%%%%%%%%%%%


\section{Measurability}

\begin{definition}[$\mu$-Measurable]
	
\end{definition}

\begin{definition}[Universally measurable]
	
\end{definition}

\begin{theorem}
	Every finite Borel measure on Polish space is regular
\end{theorem}

\begin{theorem}
	Let $B$ be an analytic subset of a polish space $X$. Then $B$ is universally measurable.
\end{theorem}

\begin{theorem}
	Let $(X,\mathcal{A})$,  $\left( Y, \mathcal{B} \right) $ be measurable spaces, that is, spaces endowed with a $\sigma$-Algebra.  
	Let $\mathcal{A}_*$ and $\mathcal{B}_*$ be the $\sigma$-Algebras of universally measurable sets.
	If  $f:X \to Y$ is  $\mathcal{A}- \mathcal{B}$-measurable, then it is also $\mathcal{A}_*-\mathcal{B}_*$-measurable
\end{theorem}


\begin{definition}[Analytic space]
	
\end{definition}

\begin{theorem}
	Let X,A be analytic meas. space, Y polish, f measurable then f(A) analytic.
\end{theorem}


%%%%%%%%%%%%%%%%%%%%%%%%%%%%%%%%%%%%%%%%%%%%%%%%%%%%%%%%%%%%%


\section{Alternative description}


\subsection{Souslin operation}
[Rogers, p.319 ?]

\subsection{Luzins construction ??}
[Sur les ensembles analytiques]


 \subsection{Projections}

%%%%%%%%%%%%%%%%%%%%%%%%%%%%%%%%%%%%%%%%%%%%%%%%%%%%%%%%%%%%%


\section{K-analytic sets??? prettyprettyplease?}
































\vspace{\fill}
%\nocite{*}
\printbibliography{}
\end{document}
